% !TEX root = hw6-answers-graded.tex
\chead{\textbf{Exercises week 6}}

\begin{document}
\flushleft\textbf{Topic: Causal Bayesian Networks, Markov Property}
\begin{enumerate}

  \item %(1 point)
  Consider the CBN with variables $V=\{X, Y, Z\}$, $J=\emptyset$, graph $G$ as given by Figure \ref{fig:one}, and Markov kernels given by
  \begin{align*}
    P(X)      & = \mathcal{N}(\mu_1, \sigma^2_1),                      \\
    P(Y|X)    & = \mathcal{N}(a\cdot X + \mu_2, \sigma^2_2),           \\
    P(Z|X, Y) & = \mathcal{N}(b\cdot Y + c\cdot X + \mu_3, \sigma^2_3)
  \end{align*}
  for given $\sigma_i^2\in\R_{>0}$ and $\mu_i\in\R$ for all $i\in\{1, 2, 3\}$.

  \begin{figure}[ht]
    \centering
    \begin{tikzpicture}[scale=1, transform shape]
      \node[ndout] (X) at (0, 1) {$X$};
      \node[ndout] (Y) at (2, 0) {$Y$};
      \node[ndout] (Z) at (0, -1) {$Z$};
      \draw[arout] (X) to (Y);
      \draw[arout] (Y) to (Z);
      \draw[arout] (X) to (Z);
    \end{tikzpicture}
    \caption{Causal graph $G$.}
    \label{fig:one}
  \end{figure}

  A Causal Bayesian Network is called \emph{faithful} with regards to its joint Markov kernel $P(X_V|\doit(X_J))$ if the inverse of the Global Markov Property holds, i.e., if for all $D, E, F \ins J \cup V$ (not-necessarily disjoint) we also have the reverse implication:
  \[ D \Perp^d_G E \given F \qquad \impliedby \qquad X_D \Indep_{P(X_V|\doit(X_J))} X_E \given X_F. \]
  Show through counterexample, by picking values $a, b, c \neq 0$, that faithfulness does not hold in general for Causal Bayesian Networks.

  \emph{Hint: Find the joint distribution of $(X, Y, Z)$, and use the fact that if a random vector has a multivariate normal distribution, then any two or more of its components that are uncorrelated (or equivalently, having zero covariance) are independent.}

  \begin{ungradedsolution}
    We have
    \begin{align*}
      \mat{X \\ Y} \sim \mathcal{N}\left( \mat{\mu_1\\ a\cdot\mu_1 + \mu_2}, \Sigma\right) \quad \text{with} \quad \Sigma := \mat{\sigma_1^2 & a\sigma_1^2 \\ a \sigma_1^2 & a^2 \sigma_1^2 + \sigma_2^2}
    \end{align*}
    and
    \begin{align*}
      P(Z|X, Y) & = \mathcal{N}\left(A\mat{X \\Y} + \mu_3, \sigma^2_3\right)
    \end{align*}
    where $A = \mat{c, b}$, and so
    \begin{align*}
      \mat{X \\ Y \\ Z} \sim \mathcal{N}\left(\mat{\mu_1\\ a\mu_1+\mu_2 \\ b(a\mu_1 + \mu_2) + c\mu_1 + \mu_3}, \mat{\Sigma & \Sigma A^\top \\ A\Sigma & \sigma_3^2 + A\Sigma A^\top}\right),
    \end{align*}
    Then we see that $\textrm{Cov}(X, Z) = (A\Sigma)_1 = (ab+c)\sigma_1^2 = 0$ whenever $ab+c=0$, so for these parameters we have $X\Indep Z$. Since $X$ is not d-separated from $Z$ in $G$, this is a faithfulness violation.
  \end{ungradedsolution}

  \item %($1 + 1$ points)
  Consider the Causal Bayesian Network over nodes $X$ and $Y$, where both variables take values in $\mathbb{R}$, where the graph is given by $X \rightarrow Y$ and the Markov kernels are defined by the densities:
  \begin{align*}
    p_X(x)   & = \mathcal{N}(x|0,1)            \\
    p_Y(y|x) & = \mathcal{N}(y|5\cdot x -3, 1)
  \end{align*}
  \begin{enumerate}
    \item Show that we can define a Causal Bayesian Network where the graph is given by $Y \rightarrow X$, that yields the same joint Markov kernel.
   % TODO 2024: rephrase: 'Define a CBN such that ...' (instead of 'Show that we can define...')
    \item What is the joint Markov kernel for both the original and the inverted Causal Bayesian Network if we perform the intervention do$(X=1)$?
    % TODO 2024: the intervention d(X=1) is not yet defined!
  \end{enumerate}

  \emph{Remark: From this exercise we see that only knowing the joint distribution of the variables of interest is not enough to be able to predict what will happen to the system if we intervene on it. Knowledge of the causal graph is essential for e.g.\ predicting the effect of a newly introduced political policy, or knowing whether eating more chocolate will increase your probability of winning a Nobel price.}

  \begin{gradedsolution}
    \begin{enumerate}
      \item \points{2} We know that if $X\sim\mathcal{N}(\mu_x,\sigma^2_x)$ and $Y\sim\mathcal{N}(a\cdot X+b, \sigma^2_y)$ are univariate normal, then their joint distribution is given by:
      \begin{align*}
        \mat{X \\ Y} &\sim \mathcal{N}\left( \mat{\mu_x\\ a\cdot\mu_x + b}, \mat{\sigma_x^2 & a\sigma_x^2 \\ a \sigma_x^2 & a^2 \sigma_x^2 + \sigma_y^2}\right),
      \end{align*}
      where in this case
      \begin{align*}
        \mat{X \\ Y} &\sim \mathcal{N}\left( \mat{0\\ -3}, \mat{1 & 5 \\ 5 & 26}\right).
      \end{align*}
      Our goal then is to define Markov kernels $P(Y)$ and $P(X|Y)$ that yield the same joint distribution. We first can ensure that the marginal distribution of $Y$ is the same by setting $P(Y) = \mathcal{N}(\mu'_y, \sigma'^2_y) = \mathcal{N}(a\cdot\mu_x + b, a^2 \sigma_x^2 + \sigma_y^2)$, thus $\mu'_y=-3$ and $\sigma'^2_y=26$. Then we define the Markov kernel $P(X|Y=y) = \mathcal{N}(a'y + b', \sigma'^2_x)$ such that the implied joint distribution remains the same. This yields the equations:
      \begin{align*}
        a\sigma_x^2 & = a' \sigma'^2_y,                 \\
        \mu_x       & = a' \mu'_y + b',                 \\
        \sigma^2_x  & = a'^2 \sigma'^2_y + \sigma'^2_x.
      \end{align*}
      The first equation immediately yields $a' = \frac{a\sigma_x^2}{\sigma'^2_y} = \frac{5}{26}$, the second equation yields $b' = \frac{15}{26}$, and the final equation yields after substitution $\sigma'^2_x = \frac{1}{26}$, giving
      \begin{align*}
        P(Y) &= \Ncal(-3, 26) \\
        P(X|Y=y) &= \Ncal(\frac{5}{26}y + \frac{15}{26}, \frac{1}{26}).
      \end{align*}
      Thus we have defined Markov kernels such that the graph $Y \rightarrow X$ yields the same joint distribution.
      \item \points{1} In both cases, after the intervention $\doit(X=1)$, $X$ is distributed according to a delta distribution with its mass point on $X=1$. In the original Causal Bayesian Network, after intervention do$(X=1)$, the marginal distribution of $Y$ changes from $\mathcal{N}(-3,26)$ to $\mathcal{N}(2,1)$. However, in the inverted network, Intervention on $X$ does not change the distribution of $Y$ and so it is still distributed as $\mathcal{N}(-3,26)$.
    \end{enumerate}
  \end{gradedsolution}

  \item %(1 point)
  Consider the generative process defined by the following densities over continuous variables $V, W, X, Y$, and probability mass function over discrete variable $Z$:
  \begin{align*}
    p_V(v)            & = \mathcal{N}(0,1),       \\
    p_W(w\given v, x) & = \text{Gamma}(v^2, x^2)  \\
    p_X(x\given u)    & = \text{Exponential}(u)   \\
    p_U(u)            & = \text{Weibull}(5,3)     \\
    p_Y(y\given v, x) & = \text{Pareto}(v^2, x^2) \\
    p_Z(z\given y)    & = \text{Poisson}(y^2)
  \end{align*}
  Show that $V \Indep Z \given Y$ and $V \Indep X$.

  \begin{ungradedsolution}
    \begin{figure}[ht]
      \centering
      \begin{tikzpicture}[scale=1, transform shape]
        \node[ndout] (V) at (0, 0) {$V$};
        \node[ndout] (W) at (2, 0) {$W$};
        \node[ndout] (X) at (4, 0) {$X$};
        \node[ndout] (U) at (6, 0) {$U$};
        \node[ndout] (Y) at (2, -2) {$Y$};
        \node[ndout] (Z) at (2, -4) {$Z$};
        \draw[arout] (V) to (W);
        \draw[arout] (X) to (W);
        \draw[arout] (U) to (X);
        \draw[arout] (V) to (Y);
        \draw[arout] (X) to (Y);
        \draw[arout] (Y) to (Z);
      \end{tikzpicture}
      \caption{Causal graph $G$ corresponding to exercise 3.}
      \label{fig:three}
    \end{figure}
    We draw a corresponding graph in figure \ref{fig:three}. From this, using the Markov property, we can immediately read off the d-separations $V \Perp^d_G Z \given Y$ and $V \Perp^d_G X$, thus the corresponding independencies are implied by the Markov property.
  \end{ungradedsolution}



\end{enumerate}
\end{document}
